% This is a figure template that can be included in the main document
% You can replace this with actual figure code or use \includegraphics

\begin{figure}[t]
    \centering
    \begin{tabular}{ll}
        \toprule
        \textbf{Feature} & \textbf{Parameter Value} \\
        \midrule
        ([],`') & 1.0 \\
        ([1],`A') & 0.5 \\
        ([1],`C') & -0.3 \\
        ([1],`G') & 0.2 \\
        ([1],`T') & -0.4 \\
        ([2],`A') & 0.1 \\
        ([1,2],`AC') & 0.7 \\
        ([1,2],`GT') & -0.5 \\
        \bottomrule
    \end{tabular}
    \caption{Example of how features are represented in GaugeFixer. Each feature is represented by a tuple containing a list of positions and a subsequence, and each parameter value corresponds to the weight for that feature. The first row represents the constant (order-0) term.}
    \label{fig:features}
\end{figure}